\documentclass[12pt]{article}
\usepackage[margin=1in]{geometry}
\usepackage{amsmath}
\usepackage{amssymb}
\usepackage{amsthm}
\usepackage{enumitem}
\usepackage{xcolor}
\usepackage{pagecolor}
\usepackage{marginnote}
\usepackage{fancyhdr}
\usepackage{bm}
\usepackage{etoolbox}

\definecolor{cream}{HTML}{faf7f1}
\pagecolor{cream}

\newcommand{\defn}[1]{\textbf{#1}}
\newcommand{\defeq}{\mathrel{\vcenter{\baselineskip0.5ex \lineskiplimit0pt
                     \hbox{\scriptsize.}\hbox{\scriptsize.}}}=}
\newcommand{\T}{\top}
\newcommand{\F}{\bot}
\newcommand{\Pow}{\mathbb{P}}
\newcommand{\N}{\mathbb{N}}

\newtheoremstyle{proofstyle}
  {}{}{\itshape}{}{\bfseries}{.}{ }{}
\theoremstyle{proofstyle}
\newtheorem*{proof*}{Proof}

\AtBeginEnvironment{proof*}{\setlength{\parskip}{0.5em}}

\newcommand{\sidenote}[1]{\marginnote{\small\itshape #1}}

\makeatletter
\renewcommand{\maketitle}{%
  \begin{flushleft}
    {\Large\bfseries CS173: Discrete Structures}\\[0.3em]
    {\Large\bfseries Problem Set}\\[0.5em]
    {\normalsize Mohammed Al Salhi}
  \end{flushleft}
  \vspace{1em}
}
\makeatother

\AtBeginDocument{\mathversion{bold}}

\begin{document}

\maketitle

\begin{enumerate}[left=0pt, itemsep=2em, parsep=1em]

\item Prove $(\forall x \in \mathbb{N})(S(x) \neq \emptyset$).

\begin{proof*}
    Let $n \in \mathbb{N}$. We proceed by contradiction. Assume $S(n) = \emptyset$.

    By definition, $S(n) = n \cup \{n\}$, so we have:
    \[
        n \cup \{n\} = \emptyset.
    \]
    But $n \in \{n\}$, so $n \in n \cup \{n\}$, which means $n \cup \{n\} \neq \emptyset$ is a 
    a contradiction.

    Therefore $S(n) \neq \emptyset$ for all $n \in \mathbb{N}$. \qed
\end{proof*}

\item Prove $(\forall n \in \mathbb{N})(n \neq 0\Rightarrow (\exists m \in \mathbb{N})(S(m) = n)$.

\begin{proof*}
    Let $n \in \mathbb{N}$ and assume $n \neq 0$.

    By the definition of natural numbers, every natural number is either $0$
    or the successor of some natural number. Since $n \neq 0$, it must be the
    case that $n = S(m)$ for some $m \in \mathbb{N}$.

    This $m$ is our witness and $S(m) = n$. \qed
\end{proof*}

\item Prove $(\forall n \in \mathbb{N})(\forall m \in \mathbb{N})(S(n + m) = S(n) + m)$.

\begin{proof*}
    Let $n \in \mathbb{N}$. We induct on $m$.

    \textbf{Base case:} $m = 0$. We need to show $S(n + 0) = S(n) + 0$.
    \[
        S(n + 0) = S(n) = S(n) + 0
    \]
    where the first equality is by definition of addition and the second is by
    definition of addition.

    \textbf{Inductive hypothesis:} Assume $S(n + k) = S(n) + k$ for some
    $k \in \mathbb{N}$.

    \textbf{Inductive step:} We need to show $S(n + S(k)) = S(n) + S(k)$.
    \[
        S(n + S(k)) = S(S(n + k)) = S(S(n) + k) = S(n) + S(k)
    \]
    where the first equality is by definition of addition, the second is by
    the inductive hypothesis, and the third is by definition of addition. \qed
\end{proof*}

\item Prove $(\forall n \in \mathbb{N})(0 \leq n)$.

\textbf{Proof by induction:}

\begin{proof*}
    We induct on $n$.

    \textbf{Base case:} $n = 0$. We need to show $0 \leq 0$, which holds
    immediately since $\leq$ is reflexive.

    \textbf{Inductive hypothesis:} Assume $0 \leq k$ for some $k \in \mathbb{N}$.

    \textbf{Inductive step:} We need to show $0 \leq S(k)$.
    By the inductive hypothesis $0 \leq k$, and by definition $k < S(k)$,
    so by transitivity $0 \leq S(k)$. \qed
\end{proof*}

\textbf{Proof by definition of $\leq$:}

\begin{proof*}
    Let $n \in \mathbb{N}$. By definition, $0 \leq n$ if and only if there
    exists $m \in \mathbb{N}$ such that $0 + m = n$. Taking $m = n$, we have
    $0 + n = n$ by definition of addition. Hence $0 \leq n$. \qed
\end{proof*}

\item Prove $(\forall n \in \mathbb{N})(n \neq 0 \Rightarrow 0 < n)$.

\begin{proof*}
    We induct on $n$.

    \textbf{Base case:} NTS: $0 \neq 0 \Rightarrow 0 < 0$.
    Assume $0 \neq 0$. We know $(\forall x)(x = x)$, so $0 = 0$ ---
    a contradiction. Thus the implication holds vacuously.

    \textbf{Inductive hypothesis:} Assume $k \neq 0 \Rightarrow 0 < k$
    for some $k \in \mathbb{N}$.

    \textbf{Inductive step:} NTS: $S(k) \neq 0 \Rightarrow 0 < S(k)$.
    Assume $S(k) \neq 0$. By definition of $<$, we have $0 < S(k)$
    if and only if $0 \leq k$, which holds by the previous theorem. \qed
\end{proof*}

\item Prove $(\forall x \in \mathbb{N})(\forall y \in \mathbb{N})(x + y = y + x)$.

\textbf{Lemma:} $(\forall n \in \mathbb{N})(\forall x \in \mathbb{N})(S(n) + x = S(n + x))$.

\begin{proof*}
    Let $n \in \mathbb{N}$. We induct on $x$.

    \textbf{Base case:} NTS: $S(n) + 0 = S(n + 0)$.
    \[
        S(n) + 0 = S(n) = S(n + 0)
    \]
    where the first equality is by definition of addition and the second is
    by definition of addition.

    \textbf{Inductive hypothesis:} Assume $S(n) + k = S(n + k)$
    for some $k \in \mathbb{N}$.

    \textbf{Inductive step:} NTS: $S(n) + S(k) = S(n + S(k))$.
    \[
        S(n) + S(k) = S(S(n) + k) = S(S(n + k)) = S(n + S(k))
    \]
    where the first equality is by definition of addition, the second is by
    the inductive hypothesis, and the third is by definition of addition. \qed
\end{proof*}

\textbf{Main proof:}

\begin{proof*}
    Let $x \in \mathbb{N}$. We induct on $y$.

    \textbf{Base case:} NTS: $x + 0 = 0 + x$.
    By definition of addition, $x + 0 = x$. By theorem, $0 + x = x$.
    Hence $x + 0 = 0 + x$.

    \textbf{Inductive hypothesis:} Assume $x + k = k + x$
    for some $k \in \mathbb{N}$.

    \textbf{Inductive step:} NTS: $x + S(k) = S(k) + x$.
    \[
        x + S(k) = S(x + k) = S(k + x) = S(k) + x
    \]
    where the first equality is by definition of addition, the second is by
    the inductive hypothesis, and the third is by the lemma. \qed
\end{proof*}

\item Prove $(\forall x \in \mathbb{N})((x \cdot 1 = x) \wedge (1 \cdot x = x))$.

\begin{proof*}
    We induct on $x$.

    \textbf{Base case:} NTS: $(0 \cdot 1 = 0) \wedge (1 \cdot 0 = 0)$.

    Consider $0 \cdot 1$:
    \[
        0 \cdot 1 = 0 \cdot S(0) = (0 \cdot 0) + 0 = 0 + 0 = 0
    \]
    where the first equality is by definition of $1 = S(0)$, the second is
    by definition of multiplication, and the third and fourth are by
    definition of addition.

    Consider $1 \cdot 0$:
    \[
        1 \cdot 0 = 0
    \]
    by definition of multiplication.

    \textbf{Inductive hypothesis:} Assume $k \cdot 1 = k$ and $1 \cdot k = k$
    for some $k \in \mathbb{N}$.

    \textbf{Inductive step:} NTS: $S(k) \cdot 1 = S(k)$ and $1 \cdot S(k) = S(k)$.

    Consider $S(k) \cdot 1$:
    \[
        S(k) \cdot 1 = S(k) \cdot S(0) = (S(k) \cdot 0) + S(k) = 0 + S(k) = S(k)
    \]
    where the first equality is by definition of $1 = S(0)$, the second is
    by definition of multiplication, the third is by definition of
    multiplication, and the fourth is by theorem.

    Consider $1 \cdot S(k)$:
    \[
        1 \cdot S(k) = (1 \cdot k) + 1 = k + 1 = S(k)
    \]
    where the first equality is by definition of multiplication, the second
    is by the inductive hypothesis, and the third is by definition of
    $S(k) = k + 1$. \qed
\end{proof*}

\item Prove $(\forall x \in \mathbb{N})((x \cdot 0 = 0) \wedge (0 \cdot x = 0))$.

\begin{proof*}
    For the first conjunct, $x \cdot 0 = 0$ holds directly by definition
    of multiplication.

    For the second conjunct, we induct on $x$ to prove
    $(\forall x \in \mathbb{N})(0 \cdot x = 0)$.

    \textbf{Base case:} NTS: $0 \cdot 0 = 0$.
    This holds by definition of multiplication.

    \textbf{Inductive hypothesis:} Let $k \in \mathbb{N}$ and assume
    $0 \cdot k = 0$.

    \textbf{Inductive step:} NTS: $0 \cdot S(k) = 0$.
    \[
        0 \cdot S(k) = (0 \cdot k) + 0 = 0 + 0 = 0
    \]
    where the first equality is by definition of multiplication, the second
    is by the inductive hypothesis, and the third is by definition of
    addition. \qed
\end{proof*}

\item Prove $(\forall x, y, z \in \mathbb{N})(x \cdot (y + z) = (x \cdot y) + (x \cdot z))$.

\begin{proof*}
    Let $x, y \in \mathbb{N}$. We induct on $z$ to prove
    $(\forall z \in \mathbb{N})(x \cdot (y + z) = (x \cdot y) + (x \cdot z))$.

    \textbf{Base case:} NTS: $x \cdot (y + 0) = (x \cdot y) + (x \cdot 0)$.

    Consider the LHS:
    \[
        x \cdot (y + 0) = x \cdot y
    \]
    by definition of addition. Consider the RHS:
    \[
        (x \cdot y) + (x \cdot 0) = (x \cdot y) + 0 = x \cdot y
    \]
    where the first equality is by definition of multiplication and the
    second is by definition of addition. Hence LHS $=$ RHS.

    \textbf{Inductive hypothesis:} Assume $x \cdot (y + k) = (x \cdot y) + (x \cdot k)$
    for some $k \in \mathbb{N}$.

    \textbf{Inductive step:} NTS: $x \cdot (y + S(k)) = (x \cdot y) + (x \cdot S(k))$.
    \[
        x \cdot (y + S(k)) = x \cdot S(y + k) = (x \cdot (y + k)) + x
        = ((x \cdot y) + (x \cdot k)) + x
    \]
    \[
        = (x \cdot y) + ((x \cdot k) + x) = (x \cdot y) + (x \cdot S(k))
    \]
    where the first equality is by definition of addition, the second is by
    definition of multiplication, the third is by the inductive hypothesis,
    the fourth is by associativity of addition, and the fifth is by
    definition of multiplication. \qed
\end{proof*}

\item Prove $(\forall x, y \in \mathbb{N})(x \cdot y = y \cdot x)$.

\textbf{Lemma:} $(\forall k \in \mathbb{N})(\forall x \in \mathbb{N})(S(k) \cdot x = (k \cdot x) + x)$.

\begin{proof*}
    Let $k \in \mathbb{N}$. We induct on $x$.

    \textbf{Base case:} NTS: $S(k) \cdot 0 = (k \cdot 0) + 0$.
    \[
        S(k) \cdot 0 = 0 = 0 + 0 = (k \cdot 0) + 0
    \]
    where the first equality is by definition of multiplication, the second
    is by definition of addition, and the third is by definition of
    multiplication.

    \textbf{Inductive hypothesis:} Assume $S(k) \cdot x = (k \cdot x) + x$
    for some $x \in \mathbb{N}$.

    \textbf{Inductive step:} NTS: $S(k) \cdot S(x) = (k \cdot S(x)) + S(x)$.
    \[
        S(k) \cdot S(x) = (S(k) \cdot x) + S(k) = ((k \cdot x) + x) + S(k)
    \]
    \[
        = (k \cdot x) + (x + S(k)) = (k \cdot x) + S(x + k)
    \]
    \[
        = (k \cdot x) + S(k + x) = (k \cdot x) + (k + S(x))
    \]
    \[
        = (k \cdot S(x)) + S(x)
    \]
    where the first equality is by definition of multiplication, the second
    is by the inductive hypothesis, the third and fourth are by associativity
    and definition of addition, the fifth is by commutativity of addition,
    the sixth is by definition of addition, and the seventh is by definition
    of multiplication. \qed
\end{proof*}

\textbf{Main proof:}

\begin{proof*}
    Let $x \in \mathbb{N}$. We induct on $y$.

    \textbf{Base case:} NTS: $x \cdot 0 = 0 \cdot x$.
    By definition of multiplication, $x \cdot 0 = 0$. By theorem,
    $0 \cdot x = 0$. Hence $x \cdot 0 = 0 \cdot x$.

    \textbf{Inductive hypothesis:} Assume $x \cdot k = k \cdot x$
    for some $k \in \mathbb{N}$.

    \textbf{Inductive step:} NTS: $x \cdot S(k) = S(k) \cdot x$.
    \[
        x \cdot S(k) = (x \cdot k) + x = (k \cdot x) + x = S(k) \cdot x
    \]
    where the first equality is by definition of multiplication, the second
    is by the inductive hypothesis, and the third is by the lemma. \qed
\end{proof*}

\item Prove $(\forall x, y, z \in \mathbb{N})(x^y \cdot x^z = x^{y+z})$.

\begin{proof*}
    Let $x, y \in \mathbb{N}$. We induct on $z$.

    \textbf{Base case:} NTS: $x^y \cdot x^0 = x^{y+0}$.
    \[
        x^y \cdot x^0 = x^y \cdot 1 = x^y = x^{y+0}
    \]
    where the first equality is by definition of exponentiation, the second
    is by definition of multiplication, and the third is by definition of
    addition.

    \textbf{Inductive hypothesis:} Assume $x^y \cdot x^k = x^{y+k}$
    for some $k \in \mathbb{N}$.

    \textbf{Inductive step:} NTS: $x^y \cdot x^{S(k)} = x^{y+S(k)}$.
    \[
        x^y \cdot x^{S(k)} = x^y \cdot (x^k \cdot x) = (x^y \cdot x^k) \cdot x
        = x^{y+k} \cdot x = x^{S(y+k)} = x^{y+S(k)}
    \]
    where the first equality is by definition of exponentiation, the second
    is by associativity of multiplication, the third is by the inductive
    hypothesis, the fourth is by definition of exponentiation, and the fifth
    is by definition of addition. \qed
\end{proof*}

\item Prove $(\forall n \in \mathbb{N})\left(6 \cdot \sum_{i=0}^{n} i^2 = n(n+1)(2n+1)\right)$.

\begin{proof*}
    We induct on $n$.

    \textbf{Base case:} NTS: $6 \cdot \sum_{i=0}^{0} i^2 = 0 \cdot 1 \cdot 1$.
    \[
        6 \cdot \sum_{i=0}^{0} i^2 = 6 \cdot 0 = 0 = 0 \cdot 1 \cdot 1
    \]
    as required.

    \textbf{Inductive hypothesis:} Assume $6 \cdot \sum_{i=0}^{k} i^2 = k(k+1)(2k+1)$
    for some $k \in \mathbb{N}$.

    \textbf{Inductive step:} NTS: $6 \cdot \sum_{i=0}^{S(k)} i^2 = S(k)(S(k)+1)(2S(k)+1)$.
    \[
        6 \cdot \sum_{i=0}^{S(k)} i^2 = 6 \cdot \left(\sum_{i=0}^{k} i^2 + S(k)^2\right)
        = 6 \cdot \sum_{i=0}^{k} i^2 + 6 \cdot S(k)^2
    \]
    \[
        = k(k+1)(2k+1) + 6(k+1)^2
    \]
    \[
        = (k+1)\left[k(2k+1) + 6(k+1)\right]
    \]
    \[
        = (k+1)(2k^2 + k + 6k + 6)
    \]
    \[
        = (k+1)(2k^2 + 7k + 6)
    \]
    \[
        = (k+1)(k+2)(2k+3)
    \]
    \[
        = S(k)(S(k)+1)(2S(k)+1)
    \]
    where the first equality splits the sum, the second distributes the $6$,
    the third applies the inductive hypothesis and uses $S(k) = k+1$, and
    the remaining steps are algebraic simplification. \qed
\end{proof*}

\end{enumerate}

\end{document}